\chapter{Lecture 7: Pass by Reference, Dynamic Memory, Operator Overloading}
\begin{center}
\textit{January 25, 2018}
\end{center}

\section{Recap: How \code{cin >> x} Works}
\begin{lstlisting}[style=cpp]
void inc(int &n) { n += 1; }
int x = 5;
inc(x);
cout << x << endl;   // prints 6
\end{lstlisting}
\code{cin >> x} is really the function call \code{operator>>(cin, x)}, where
\begin{lstlisting}[style=cpp]
std::istream &operator>>(std::istream &in, int &data);
\end{lstlisting}
Both parameters are references: \code{data} so that the read can change
\code{x}, and \code{in} because streams cannot be copied (and so that we return
the same stream, enabling cascading).

\section{Pass by Value vs.\ Pass by Reference}
\begin{lstlisting}[style=cpp]
struct ReallyBig { /* lots of fields */ };

void f(ReallyBig rb) { ... }
ReallyBig b;
f(b);       // a copy of b is created for f: expensive
\end{lstlisting}
To avoid the copy, in C we would pass a pointer to \code{b}. That works in C++
too, but references are nicer:
\begin{lstlisting}[style=cpp]
void g(ReallyBig &rb) { ... }        // rb is another name for b
g(b);                                // no copy; but g could change b

void h(const ReallyBig &rb) { ... }  // no copy, and
h(b);                                // h cannot change b
\end{lstlisting}
\begin{center}
\begin{tikzpicture}[>=Stealth,font=\small,big/.style={draw,fill=blue!5,minimum width=2.4cm,minimum height=1cm}]
  \node[big] (b) at (0,0) {\code{b}};
  \node[font=\scriptsize,anchor=south] at (b.north) {\code{main}};
  \node[big] (c) at (0,-1.8) {\code{rb} (a copy)};
  \node[font=\scriptsize,anchor=north] at (c.south) {\code{f(ReallyBig rb)}};
  \draw[->,dashed] (b) -- node[right,font=\scriptsize]{copied} (c);
  \begin{scope}[xshift=6cm]
  \node[big] (b2) at (0,0) {\code{b}};
  \node[font=\scriptsize,anchor=south] at (b2.north) {\code{main}};
  \node[draw,minimum width=1.2cm,minimum height=0.7cm] (r) at (0,-1.8) {\code{rb}};
  \node[font=\scriptsize,anchor=north] at (r.south) {\code{h(const ReallyBig \&rb)}};
  \draw[->] (r) -- node[right,font=\scriptsize]{another name for \code{b}} (b2);
  \end{scope}
\end{tikzpicture}
\end{center}
People often call \code{rb} a ``const reference''. Technically the correct
reading is that \code{rb} is a \emph{reference to const} (a reference itself can
never be re-bound anyway).

\begin{important}[Advice]
Whenever possible, prefer to pass by reference to const for anything bigger
than an \code{int}, unless the function needs to change the argument (pass by
reference) or needs its own copy anyway (pass by value).
\end{important}

Another benefit of passing by reference to const:
\begin{lstlisting}[style=cpp]
void foo(int &x) { ... }
foo(5);      // illegal: if foo did x = 10, that would mean 5 = 10
foo(y + y);  // illegal: y + y is not an lvalue

void bar(const int &x) { ... }
bar(5);      // OK
bar(y + y);  // OK
\end{lstlisting}
C++ allows a reference to const to be initialized from a value that is not an
lvalue: the compiler creates a temporary to hold the value and binds the
reference to it. Since \code{bar} promises not to change \code{x}, nobody can
notice.

\section{Dynamic Memory Allocation (Heap Allocation)}
In C:
\begin{lstlisting}[style=c]
int *p = malloc(n * sizeof(int));
...
free(p);
\end{lstlisting}
This works in C++, but in CS 246 we do not use it. In C++ we use \code{new} and
\code{delete}. \code{new} is type-aware: it knows the size of the type, and (as
we'll see) it calls the constructor.
\begin{lstlisting}[style=cpp]
struct Node {
    int data;
    Node *next;
};

Node *np = new Node;   // no need to say how much space is needed
...
delete np;             // free memory you heap-allocated earlier
\end{lstlisting}
It is always safe to \code{delete nullptr} (it does nothing).

\begin{lstlisting}[style=cpp]
Node n;
Node *np = new Node;
\end{lstlisting}
\begin{center}
\begin{tikzpicture}[box/.style={draw, minimum width=16mm, minimum height=7mm}]
\node[font=\bfseries] at (0,1.2) {Stack};
\node[font=\bfseries] at (5,1.2) {Heap};
\node[box] (n) at (0,0.3) {\code{n}: Node};
\node[box] (np) at (0,-0.6) {\code{np}};
\node[box] (h) at (5,-0.6) {Node};
\draw[-{Stealth}] (np.east) -- (h.west);
\end{tikzpicture}
\end{center}
When \code{n} and \code{np} go out of scope, their stack memory is reclaimed.
The heap Node continues to live on.

\begin{important}[Memory leaks]
Forgetting to deallocate heap memory causes a \textbf{memory leak}. Memory leaks
are bad, and Marmoset tests run under \code{valgrind}: leaking even one byte gives
you 0 on that test.
\end{important}

\subsection{Heap Arrays}
\begin{lstlisting}[style=cpp]
Node *narr = new Node[10];
...
delete [] narr;   // arrays must be freed with delete []
\end{lstlisting}
Memory allocated with \code{new} must be freed with \code{delete}; memory
allocated with \code{new []} must be freed with \code{delete []}. Mixing them is
undefined behaviour. (Why \code{[]} is needed is discussed at the end of CS 241:
the runtime has to know how many destructors to run.)

\section{Returning Values}
\begin{lstlisting}[style=cpp]
Node getMeANode() {
    Node n;
    return n;
}
\end{lstlisting}
This returns by value, so (in principle) \code{n} is copied into the caller.
That can be expensive --- we come back to this later in the course (move
semantics and copy elision).

\begin{lstlisting}[style=cpp]
// ------- DON'T DO THIS -------
Node *getMeANode() {
    Node n;
    return &n;     // n dies when the function returns!
}
\end{lstlisting}
The returned pointer is a \textbf{dangling pointer}: it points to memory that is
no longer yours to access.

\begin{center}
\begin{tikzpicture}[>=Stealth,font=\small,cell/.style={draw,minimum width=3.2cm,minimum height=0.7cm}]
  \node[cell,fill=blue!5] (m) at (0,0) {caller: \code{Node *p}};
  \node[cell,fill=gray!15,draw=gray,text=gray,below=0pt of m] (dead) {\code{n} (frame popped)};
  \node[font=\scriptsize,anchor=west,gray] at (dead.east) {\ was \code{getMeANode}'s frame};
  \draw[->,red!60!black,thick] (m.west) to[out=180,in=180,looseness=1.6] (dead.west);
  \node[font=\scriptsize,red!60!black,anchor=east] at (-2.3,-0.4) {dangling};
\end{tikzpicture}
\end{center}

\begin{important}[Rule]
Never return a pointer or a reference to stack-allocated (local) memory.
\end{important}

Returning a pointer to heap memory is fine (but now the caller is responsible for
deleting it):
\begin{lstlisting}[style=cpp]
Node *getMeANode() {
    Node *np = new Node;
    return np;
}
\end{lstlisting}

\section{Operator Overloading}
C++ lets us give meaning to C++ operators for types we create.
\begin{lstlisting}[style=cpp]
struct Vec {
    int x, y;
};
\end{lstlisting}
In C we would write a function \code{Vec add(Vec v1, Vec v2);}. In C++:
\begin{lstlisting}[style=cpp]
Vec operator+(const Vec &v1, const Vec &v2) {
    Vec v{v1.x + v2.x, v1.y + v2.y};
    return v;
}

Vec operator*(const Vec &v, int k) {
    return {v.x * k, v.y * k};   // the compiler knows to build a Vec
}

Vec operator*(int k, const Vec &v) {
    return v * k;                // so k * v also works
}

Vec v1{1, 2}, v2{3, 4};
Vec v3 = v1 + v2 * 2;            // {7, 10}: usual precedence applies
\end{lstlisting}
Each ordering of operands needs its own function: \code{v * k} and \code{k * v}
are different calls.

\subsection{Overloading the I/O Operators}
\begin{lstlisting}[style=cpp]
struct Grade {
    int theGrade;
};

std::ostream &operator<<(std::ostream &out, const Grade &g) {
    out << g.theGrade << '%';
    return out;
}

std::istream &operator>>(std::istream &in, Grade &g) {
    in >> g.theGrade;
    if (g.theGrade < 0) g.theGrade = 0;
    else if (g.theGrade > 100) g.theGrade = 100;
    return in;
}

Grade g;
cin >> g;               // reads an int, clamped to [0, 100]
cout << g << endl;      // e.g. 85%
\end{lstlisting}
Returning the stream by reference is what makes \code{cout << g << endl} work.

\section{The Preprocessor}
The preprocessor is a program that runs \emph{before} the compiler sees your
code; it transforms your source text. Lines starting with \code{\#} are
\textbf{preprocessor directives}.
\begin{itemize}
  \item \code{\#include <file>}: copy and paste the contents of the header right
        here. Angle brackets are for standard library headers.
  \item \code{\#include "file.h"}: same, for your own headers (searched for in
        the current directory first).
\end{itemize}
To see the output of the preprocessor:
\begin{lstlisting}[style=bash]
g++ -E -P file.cc
\end{lstlisting}

\subsection{The Define Directive}
\code{\#define VAR VALUE} creates a preprocessor variable; afterwards every
occurrence of the token \code{VAR} is replaced by \code{VALUE}.
\begin{lstlisting}[style=cpp]
#define MAX 10
...
int array[MAX];   // the compiler sees: int array[10];
\end{lstlisting}
In C++ a \code{const int MAX = 10;} is usually preferable (it has a type and
obeys scope). More uses of \code{\#define} next lecture.

\section*{Summary}
\begin{itemize}
  \item Pass anything bigger than an \code{int} by reference to const; a
        reference to const can bind to temporaries.
  \item \code{new}/\code{delete} and \code{new[]}/\code{delete[]}; leaks are
        failures in this course.
  \item Never return pointers/references to locals.
  \item Overload operators by writing \code{operator+}, \code{operator<<}, etc.;
        I/O operators take and return the stream by reference.
  \item The preprocessor handles \code{\#include} and \code{\#define} before
        compilation.
\end{itemize}
